% Compile with XeLaTeX.
% 測度論最小模組：為目前的 L^2 / L^p / Banach 空間學習服務
% Preamble style follows the user's previous study documents.

\documentclass[11pt,a4paper,twoside]{article}

\usepackage{fontspec}
\usepackage{xeCJK}

\IfFontExistsTF{Times New Roman}
  {\setmainfont{Times New Roman}}
  {\setmainfont{Latin Modern Roman}}

\ifdefined\setCJKmainfont
  \IfFontExistsTF{標楷體}
    {\setCJKmainfont{標楷體}}
    {\setCJKmainfont{標楷體}}
\fi

\usepackage{amsmath,amssymb,mathtools,empheq}
\usepackage[most]{tcolorbox}
\usepackage{tikz,wrapfig}
\usepackage[a4paper,inner=2.5cm,outer=2cm,top=2.5cm,bottom=2.5cm]{geometry}
\usepackage{setspace}
\usepackage{fancyhdr}
\usepackage[hang,flushmargin]{footmisc}

\setstretch{1.25}
\setlength{\parindent}{2em}
\setlength{\parskip}{5pt}
\setlength{\abovedisplayskip}{6pt}
\setlength{\belowdisplayskip}{6pt}
\setlength{\abovedisplayshortskip}{4pt}
\setlength{\belowdisplayshortskip}{4pt}

\pagestyle{fancy}
\fancyhf{}
\renewcommand{\headrulewidth}{0pt}
\fancyfoot[RO,LE]{\small $\boldsymbol{\cdot}\;\thepage\;\boldsymbol{\cdot}$}

\newif\ifstudynotes
\studynotestrue

\newtcolorbox{studyNote}[1][]{
  enhanced,
  breakable,
  colback=gray!4,
  colframe=black!35,
  boxrule=0.5pt,
  arc=1mm,
  left=2mm,
  right=2mm,
  top=1.5mm,
  bottom=1.5mm,
  before skip=6pt,
  after skip=8pt,
  fonttitle=\bfseries,
  title={複習註解},
  #1
}

\newtcolorbox{importantNote}[1][]{
  enhanced,
  breakable,
  colback=gray!8,
  colframe=black!65,
  boxrule=0.8pt,
  arc=1mm,
  left=2mm,
  right=2mm,
  top=1.5mm,
  bottom=1.5mm,
  before skip=7pt,
  after skip=9pt,
  fonttitle=\bfseries,
  title={重要觀念},
  #1
}

\newcommand{\StudyNote}[2]{%
  \ifstudynotes
    \begin{studyNote}[title={#1}]
      #2
    \end{studyNote}
  \fi
}

\newcommand{\ImportantNote}[2]{%
  \ifstudynotes
    \begin{importantNote}[title={#1}]
      #2
    \end{importantNote}
  \fi
}

\begin{document}

\begin{center}
    {\LARGE\bfseries 測度論最小學習計劃}\\[0.4em]
    {\large 為 $L^2$、$L^p$、Banach 空間與變分法服務}
\end{center}

\par\vspace{0.4em}\hrule\vspace{1em}

\ImportantNote{學習目標}{
    目前不需要完整學完測度論。這個模組只補足理解
    $L^2$、$L^p$、幾乎處處相等、完備化與函數收斂所需的最低限度知識。

    建議安排為四次學習，每次約 60 分鐘；完成後立刻回頭重讀例 3。
}

\section*{第 1 次：測度與零測集}

先把「測度」理解成集合大小的推廣。例如區間 $[a,b]$ 的長度為
\[
    m([a,b])=b-a.
\]

目前最重要的是單點集合的測度：
\[
    m(\{x_0\})=0.
\]

有限個點，以及可數集合，也可能具有測度零。這直接解釋了例 3 中為什麼只改變
$\overline{\varphi}(0)$ 的值，不會改變 $L^2$ 距離。

\ImportantNote{本節核心}{
    若兩個函數只在一個零測集上不同，則這些差異不會影響 Lebesgue 積分。

    特別地，若 $f$ 與 $g$ 只有在 $x=0$ 不同，則
    \[
        \int |f-g|^2\,dx=0.
    \]
}

\StudyNote{練習}{
    判斷下列集合哪些測度為零：
    \[
        \{0\},\qquad
        \{1,2,3\},\qquad
        \mathbb{Q}\cap[0,1],\qquad
        [0,1].
    \]
}

\section*{第 2 次：幾乎處處相等（a.e.）}

符號
\[
    f=g\quad \text{a.e.}
\]
表示 $f$ 與 $g$ 並不一定在每一點都相等，而是它們不同的集合
\[
    \{x:f(x)\neq g(x)\}
\]
具有測度零。

例如令
\[
    f(x)=
    \begin{cases}
        1, & x\neq0, \\
        0, & x=0,
    \end{cases}
    \qquad
    g(x)=1.
\]

雖然 $f(0)\neq g(0)$，但仍有 $f=g$ a.e.

\ImportantNote{$L^p$ 空間真正的元素}{
    在 $L^p$ 空間中，只在零測集上不同的函數會被視為同一個元素。

    因此更精確地說，$L^p$ 的元素不是單一函數，而是「幾乎處處相等」所形成的等價類。
}

完成本節後，回頭重讀例 3 的段落 40--41，特別注意
$\overline{\varphi}(0)$ 為什麼可以任意指定。

\section*{第 3 次：Lebesgue 積分的基本直覺}

這一階段不學 Lebesgue 積分的完整建構，只建立使用上的直覺。

Riemann 積分主要從定義域出發切割區間；Lebesgue 積分則更適合按照函數值與對應集合來組織積分，因此在處理函數序列與極限時更自然。

目前只需要認識三個關鍵詞：

\begin{center}
    可測函數 \qquad Lebesgue 可積 \qquad $p$ 次方可積
\end{center}

對 $p\ge1$，$L^p(\Omega)$ 的核心條件是
\[
    \int_{\Omega}|f(x)|^p\,dx<\infty.
\]

若把「幾乎處處相等」也納入，則概念上可以寫成
\[
    L^p(\Omega)
    =
    \left\{
    f:\int_{\Omega}|f|^p\,dx<\infty
    \right\}\big/\!\sim,
\]
其中
\[
    f\sim g
    \quad\Longleftrightarrow\quad
    f=g\ \text{a.e.}
\]

\section*{第 4 次：三種收斂不要混在一起}

\subsection*{1. 逐點收斂}

固定 $x$，再讓 $n\to\infty$：
\[
    f_n(x)\to f(x).
\]

\subsection*{2. 幾乎處處收斂}

允許在某個零測集上收斂失敗：
\[
    f_n(x)\to f(x)\quad \text{a.e.}
\]

\subsection*{3. $L^p$ 收斂}

以 $L^p$ 範數衡量整體誤差：
\[
    \|f_n-f\|_p\to0.
\]

在 $p=2$ 時，
\[
    \|f_n-f\|_2^2
    =
    \int |f_n-f|^2\,dx
    \to0.
\]

\ImportantNote{最容易混淆的地方}{
    $L^2$ 收斂並不等於「每一點都逐點收斂」。

    例 3 的核心正是在比較「函數序列屬於 $C([-1,1])$」與
    「其 $L^2$ 極限是否仍屬於 $C([-1,1])$」。
}

\section*{完成後：回頭重讀例 3}

完成上述四次學習後，重新整理例 3：
\[
    \varphi_n\in C([-1,1]),
    \qquad
    \{\varphi_n\}\text{ 是 }L^2\text{-Cauchy 序列}.
\]

其 $L^2$ 極限可表示為
\[
    \varphi(x)=
    \begin{cases}
        -1, & x<0, \\
        1,  & x>0,
    \end{cases}
\]
而 $x=0$ 的值在 $L^2$ 意義下可以任意指定，因為單點集合的測度為零。

因此這個極限不屬於 $C([-1,1])$，卻屬於 $L^2([-1,1])$，所以
$C([-1,1])$ 在 $L^2$ 範數下不是完備空間，而其完備化為
\[
    L^2([-1,1]).
\]

\section*{目前暫時不用深入}

為了維持學習效率，現階段先不深入：

\begin{itemize}
    \item $\sigma$-代數的完整抽象理論；
    \item 外測度與 Carath\'eodory 可測性的完整建構；
    \item Radon--Nikodym 定理；
    \item Fubini 與 Tonelli 定理的完整證明；
    \item 單調收斂定理與控制收斂定理的正式證明。
\end{itemize}

等之後進入 Sobolev 空間、弱收斂或更深入的 PDE，再進行第二階段測度論。

\end{document}
